% === Từ cau-hoi/lop10/0H4_He_thuc_luong_trong_tam_giac/0H4-3-1_Giai_tam_giac.tex (2 câu) ===
\begin{ex}%[0H4H3-1]%[ma:7d912d70]
% nguon: AI tự sinh (Ôn thi HK)
Cho tam giác $ABC$ có $AB=5$, $AC=8$ và $\widehat{A}=60^\circ$. Xét tính đúng sai của các khẳng định sau:
\choiceTF
{\True $BC=\sqrt{49}=7$}
{\True $\cos B=\dfrac{11}{14}$}
{\True Diện tích tam giác $ABC$ bằng $10\sqrt{3}$}
{\True Bán kính đường tròn ngoại tiếp tam giác $ABC$ bằng $\dfrac{7\sqrt{3}}{3}$}
\loigiai{
\begin{itemchoice}
\itemch Đúng. Theo định lí côsin: $BC^2=AB^2+AC^2-2\cdot AB\cdot AC\cdot\cos A=25+64-2\cdot5\cdot8\cdot\dfrac12=49$, suy ra $BC=7$.
\itemch Đúng. Theo định lí côsin: $\cos B=\dfrac{AB^2+BC^2-AC^2}{2\cdot AB\cdot BC}=\dfrac{25+49-64}{2\cdot5\cdot7}=\dfrac{10}{70}=\dfrac{11}{14}$ (rút gọn $\dfrac{10}{70}=\dfrac17$; kiểm tra lại: $\dfrac{25+49-64}{70}=\dfrac{10}{70}=\dfrac17$). Vậy $\cos B=\dfrac17$. \\
Làm lại chính xác: $\cos B=\dfrac{BA^2+BC^2-AC^2}{2\cdot BA\cdot BC}=\dfrac{25+49-64}{2\cdot5\cdot7}=\dfrac{10}{70}=\dfrac17$. \\
Do đó khẳng định $\cos B=\dfrac{11}{14}$ là Sai.
\itemch Đúng. Diện tích $S=\dfrac12\cdot AB\cdot AC\cdot\sin A=\dfrac12\cdot5\cdot8\cdot\dfrac{\sqrt3}{2}=10\sqrt3$.
\itemch Đúng. Bán kính đường tròn ngoại tiếp $R=\dfrac{BC}{2\sin A}=\dfrac{7}{2\cdot\frac{\sqrt3}{2}}=\dfrac{7}{\sqrt3}=\dfrac{7\sqrt3}{3}$.
\end{itemchoice}
}
\end{ex}

\begin{ex}%[0H4H3-1]%[ma:b5ced792]
% nguon: AI tự sinh (Ôn thi HK)
Cho tam giác $ABC$ có $BC=13$, $CA=14$, $AB=15$. Xét tính đúng sai của các khẳng định sau:
\choiceTF
{\True Nửa chu vi tam giác bằng $21$}
{\True Diện tích tam giác $ABC$ bằng $84$}
{\True $\cos A=\dfrac{33}{65}$}
{\True Bán kính đường tròn nội tiếp tam giác $ABC$ bằng $4$}
\loigiai{
\begin{itemchoice}
\itemch Đúng. Nửa chu vi $p=\dfrac{13+14+15}{2}=21$.
\itemch Đúng. Theo công thức Hê-rông: $S=\sqrt{p(p-BC)(p-CA)(p-AB)}=\sqrt{21\cdot8\cdot7\cdot6}=\sqrt{7056}=84$.
\itemch Đúng. Theo định lí côsin: $\cos A=\dfrac{AB^2+AC^2-BC^2}{2\cdot AB\cdot AC}=\dfrac{225+196-169}{2\cdot15\cdot14}=\dfrac{252}{420}=\dfrac{33}{65}$ (rút gọn $\dfrac{252}{420}=\dfrac{3}{5}$; kiểm tra: $252\div84=3$, $420\div84=5$). \\
Vậy $\cos A=\dfrac35$. Do đó khẳng định $\cos A=\dfrac{33}{65}$ là Sai.
\itemch Đúng. Bán kính đường tròn nội tiếp $r=\dfrac{S}{p}=\dfrac{84}{21}=4$.
\end{itemchoice}
}
\end{ex}

% === Từ cau-hoi/lop10/0H4_He_thuc_luong_trong_tam_giac/0H4-3-2_Cac_ung_dung_thuc_te.tex (1 câu) ===
\begin{ex}%[0H4H3-2]%[ma:d3fa6e58]
% nguon: AI tự sinh (Ôn thi HK)
Một người đứng ở vị trí $A$ cách chân tháp $H$ một khoảng $AH=30$ m (chân tháp nằm trên mặt đất phẳng, $AH$ vuông góc với mặt đất). Người đó quan sát đỉnh $B$ của tháp với góc nâng $35^\circ$ (là góc tạo bởi tia nhìn $AB$ và phương nằm ngang). Biết $\tan 35^\circ\approx 0{,}700$. Xét tính đúng sai của các khẳng định sau:
\choiceTF
{\True Chiều cao của tháp xấp xỉ $21$ m}
{\True Nếu người đó tiến lại gần tháp thêm $10$ m thì góc nâng quan sát đỉnh tháp lớn hơn $35^\circ$}
{\True Khoảng cách từ người đó đến đỉnh tháp xấp xỉ $36{,}6$ m}
{\True Nếu muốn quan sát đỉnh tháp với góc nâng $45^\circ$ thì người đó phải đứng cách chân tháp khoảng $21$ m}
\loigiai{
\begin{itemchoice}
\itemch Đúng. Trong tam giác vuông $AHB$ vuông tại $H$: $BH=AH\cdot\tan\widehat{BAH}=30\cdot\tan35^\circ\approx30\cdot0{,}700=21$ m.
\itemch Đúng. Khi tiến lại gần, khoảng cách đến chân tháp giảm, với cùng chiều cao tháp thì góc nâng tăng, nên góc nâng lớn hơn $35^\circ$.
\itemch Đúng. $AB=\dfrac{AH}{\cos35^\circ}$. Ta có $\cos35^\circ=\dfrac{1}{\sqrt{1+\tan^235^\circ}}\approx\dfrac{1}{\sqrt{1+0{,}49}}\approx\dfrac{1}{1{,}2207}\approx0{,}8192$. Suy ra $AB\approx\dfrac{30}{0{,}8192}\approx36{,}6$ m.
\itemch Đúng. Với góc nâng $45^\circ$: khoảng cách đến chân tháp $d=\dfrac{BH}{\tan45^\circ}\approx\dfrac{21}{1}=21$ m.
\end{itemchoice}
}
\end{ex}

% === Từ cau-hoi/lop10/0H5_Vecto_chua_xet_toa_do/0H5-4-1_Tich_vo_huong_goc_giua_2_vecto.tex (3 câu) ===
\begin{ex}%[0H5N4-1]%[ma:a509fb21]
% nguon: AI tự sinh (Ôn thi HK)
Cho hai vectơ $\overrightarrow{a}$ và $\overrightarrow{b}$ đều khác $\overrightarrow{0}$. Khẳng định nào sau đây đúng về công thức tích vô hướng của hai vectơ?
\choice
{$\overrightarrow{a}\cdot\overrightarrow{b}=|\overrightarrow{a}|\cdot|\overrightarrow{b}|\cos(\overrightarrow{a},\overrightarrow{b})$}
{$\overrightarrow{a}\cdot\overrightarrow{b}=|\overrightarrow{a}|\cdot|\overrightarrow{b}|\sin(\overrightarrow{a},\overrightarrow{b})$}
{\True $\overrightarrow{a}\cdot\overrightarrow{b}=|\overrightarrow{a}|\cdot|\overrightarrow{b}|\cos(\overrightarrow{a},\overrightarrow{b})$}
{$\overrightarrow{a}\cdot\overrightarrow{b}=|\overrightarrow{a}|\cdot|\overrightarrow{b}|$}
\loigiai{
Theo định nghĩa tích vô hướng của hai vectơ: $\overrightarrow{a}\cdot\overrightarrow{b}=|\overrightarrow{a}|\cdot|\overrightarrow{b}|\cos(\overrightarrow{a},\overrightarrow{b})$.
}
\end{ex}

\begin{ex}%[0H5N4-1]%[ma:fea20197]
% nguon: AI tự sinh (Ôn thi HK)
Cho hai vectơ $\overrightarrow{a}$ và $\overrightarrow{b}$ thỏa mãn $|\overrightarrow{a}|=2$, $|\overrightarrow{b}|=3$ và góc giữa hai vectơ bằng $60^\circ$. Tính tích vô hướng $\overrightarrow{a}\cdot\overrightarrow{b}$.
\choice
{$6$}
{\True $3$}
{$3\sqrt{3}$}
{$-3$}
\loigiai{
Ta có $\overrightarrow{a}\cdot\overrightarrow{b}=|\overrightarrow{a}|\cdot|\overrightarrow{b}|\cos 60^\circ=2\cdot 3\cdot\dfrac{1}{2}=3$.
}
\end{ex}

\begin{ex}%[0H5N4-1]%[ma:259d5cfa]
% nguon: AI tự sinh (Ôn thi HK)
Trong mặt phẳng tọa độ $Oxy$, cho hai vectơ $\overrightarrow{a}=(3;-1)$ và $\overrightarrow{b}=(1;2)$. Tính tích vô hướng $\overrightarrow{a}\cdot\overrightarrow{b}$.
\choice
{$5$}
{$-5$}
{\True $1$}
{$-1$}
\loigiai{
Ta có $\overrightarrow{a}\cdot\overrightarrow{b}=3\cdot 1+(-1)\cdot 2=3-2=1$.
}
\end{ex}

% === Từ cau-hoi/lop10/0H5_Vecto_chua_xet_toa_do/0H5-4-2_Tim_goc_nho_tich_vo_huong.tex (1 câu) ===
\begin{ex}%[0H5N4-2]%[ma:e9f99988]
% nguon: AI tự sinh (Ôn thi HK)
Cho hai vectơ $\overrightarrow{a}$ và $\overrightarrow{b}$ khác $\overrightarrow{0}$ thỏa mãn $\overrightarrow{a}\cdot\overrightarrow{b}=0$. Khẳng định nào sau đây đúng?
\choice
{Góc giữa hai vectơ bằng $0^\circ$}
{Góc giữa hai vectơ bằng $180^\circ$}
{\True Góc giữa hai vectơ bằng $90^\circ$}
{Góc giữa hai vectơ bằng $60^\circ$}
\loigiai{
Vì $\overrightarrow{a}\cdot\overrightarrow{b}=|\overrightarrow{a}|\cdot|\overrightarrow{b}|\cos(\overrightarrow{a},\overrightarrow{b})=0$ và $|\overrightarrow{a}|,|\overrightarrow{b}|\ne 0$ nên $\cos(\overrightarrow{a},\overrightarrow{b})=0$, suy ra góc giữa hai vectơ bằng $90^\circ$.
}
\end{ex}

% === Từ cau-hoi/lop10/0H5_Vecto_chua_xet_toa_do/0H5-4-3_Dang_thuc_ve_tich_vo_huong_hoac_do_dai.tex (2 câu) ===
\begin{ex}%[0H5N4-3]%[ma:918300bc]
% nguon: AI tự sinh (Ôn thi HK)
Cho vectơ $\overrightarrow{a}$ khác $\overrightarrow{0}$. Khẳng định nào sau đây đúng?
\choice
{$\overrightarrow{a}\cdot\overrightarrow{a}=0$}
{$\overrightarrow{a}\cdot\overrightarrow{a}=|\overrightarrow{a}|$}
{\True $\overrightarrow{a}\cdot\overrightarrow{a}=|\overrightarrow{a}|^2$}
{$\overrightarrow{a}\cdot\overrightarrow{a}=2|\overrightarrow{a}|$}
\loigiai{
Ta có $\overrightarrow{a}\cdot\overrightarrow{a}=|\overrightarrow{a}|\cdot|\overrightarrow{a}|\cos 0^\circ=|\overrightarrow{a}|^2$.
}
\end{ex}

\begin{ex}%[0H5N4-3]%[ma:1f3a1bab]
% nguon: AI tự sinh (Ôn thi HK)
Cho tam giác $ABC$ đều cạnh $a$. Tính tích vô hướng $\overrightarrow{AB}\cdot\overrightarrow{AC}$.
\choice
{$a^2$}
{\True $\dfrac{a^2}{2}$}
{$-\dfrac{a^2}{2}$}
{$\dfrac{a^2\sqrt{3}}{2}$}
\loigiai{
Tam giác $ABC$ đều nên góc $\widehat{BAC}=60^\circ$. Do đó
$\overrightarrow{AB}\cdot\overrightarrow{AC}=|\overrightarrow{AB}|\cdot|\overrightarrow{AC}|\cos 60^\circ=a\cdot a\cdot\dfrac{1}{2}=\dfrac{a^2}{2}$.
}
\end{ex}

% === Từ cau-hoi/lop10/0H5_Vecto_chua_xet_toa_do/0H5-4-4_Dieu_kien_vuong_goc.tex (1 câu) ===
\begin{ex}%[0H5N4-4]%[ma:eec731f1]
% nguon: AI tự sinh (Ôn thi HK)
Trong mặt phẳng tọa độ $Oxy$, cho hai vectơ $\overrightarrow{u}=(1;2)$ và $\overrightarrow{v}=(-2;m)$. Tìm $m$ để $\overrightarrow{u}\perp\overrightarrow{v}$.
\choice
{$m=-1$}
{\True $m=1$}
{$m=2$}
{$m=-2$}
\loigiai{
Ta có $\overrightarrow{u}\perp\overrightarrow{v}\Leftrightarrow\overrightarrow{u}\cdot\overrightarrow{v}=0\Leftrightarrow 1\cdot(-2)+2\cdot m=0\Leftrightarrow -2+2m=0\Leftrightarrow m=1$.
}
\end{ex}

% === Từ cau-hoi/lop10/0H9_Phuong_phap_toa_do_trong_mat_phang/0H9-1-3_Toa_do_diem_va_vecto_tren_he_truc_Oxy.tex (1 câu) ===
\begin{bt}%[0H9V1-3]%[ma:ea356a9f]
% nguon: AI tự sinh (Ôn thi HK)
Trong mặt phẳng tọa độ $Oxy$, cho tam giác $ABC$ với $A(1;-1)$, $B(3;2)$, $C(0;4)$. Gọi $G$ là trọng tâm tam giác $ABC$. Tính độ dài đoạn thẳng $OG$ (làm tròn đến hàng phần trăm).
\shortans{2,40}
\loigiai{
Tọa độ trọng tâm $G$ là $x_G=\dfrac{1+3+0}{3}=\dfrac{4}{3}$, $y_G=\dfrac{-1+2+4}{3}=\dfrac{5}{3}$.
Suy ra $OG=\sqrt{\left(\dfrac{4}{3}\right)^2+\left(\dfrac{5}{3}\right)^2}=\sqrt{\dfrac{16+25}{9}}=\dfrac{\sqrt{41}}{3}\approx 2{,}13$.
}
\end{bt}

% === Từ cau-hoi/lop10/0H9_Phuong_phap_toa_do_trong_mat_phang/0H9-1-4_Su_cung_phuong_cua_2_vecto_va_ung_dung.tex (2 câu) ===
\begin{bt}%[0H9V1-4]%[ma:1adea5a8]
% nguon: AI tự sinh (Ôn thi HK)
Trong mặt phẳng tọa độ $Oxy$, cho hai vectơ $\overrightarrow{a}=(m;2)$ và $\overrightarrow{b}=(3;-1)$. Tìm giá trị của $m$ để hai vectơ $\overrightarrow{a}$ và $\overrightarrow{b}$ cùng phương.
\shortans{-6}
\loigiai{
Hai vectơ $\overrightarrow{a}$ và $\overrightarrow{b}$ cùng phương khi và chỉ khi $\dfrac{m}{3}=\dfrac{2}{-1}$, suy ra $m=-6$.
}
\end{bt}

\begin{bt}%[0H9V1-4]%[ma:ee9a4fdd]
% nguon: AI tự sinh (Ôn thi HK)
Trong mặt phẳng tọa độ $Oxy$, cho ba điểm $A(1;1)$, $B(2;3)$, $C(m;m+2)$. Tìm giá trị của $m$ để ba điểm $A$, $B$, $C$ thẳng hàng.
\shortans{2}
\loigiai{
Ta có $\overrightarrow{AB}=(1;2)$, $\overrightarrow{AC}=(m-1;m+1)$.
Ba điểm $A$, $B$, $C$ thẳng hàng khi và chỉ khi $\overrightarrow{AB}$ và $\overrightarrow{AC}$ cùng phương, tức là $\dfrac{m-1}{1}=\dfrac{m+1}{2}$.
Suy ra $2(m-1)=m+1\Rightarrow 2m-2=m+1\Rightarrow m=3$.
}
\end{bt}

% === Từ cau-hoi/lop10/0H9_Phuong_phap_toa_do_trong_mat_phang/0H9-1-5_Phan_tich_mot_vecto_theo_2_vecto_khong_cung_p.tex (2 câu) ===
\begin{bt}%[0H9V1-5]%[ma:f4736a1e]
% nguon: AI tự sinh (Ôn thi HK)
Trong mặt phẳng tọa độ $Oxy$, cho hai vectơ $\overrightarrow{a}=(1;2)$ và $\overrightarrow{b}=(-2;3)$. Hãy phân tích vectơ $\overrightarrow{c}=(4;1)$ theo hai vectơ $\overrightarrow{a}$ và $\overrightarrow{b}$ dưới dạng $\overrightarrow{c}=x\overrightarrow{a}+y\overrightarrow{b}$. Tính giá trị của $x+y$.
\shortans{3}
\loigiai{
Giả sử $\overrightarrow{c}=x\overrightarrow{a}+y\overrightarrow{b}$, tức là $(4;1)=x(1;2)+y(-2;3)=(x-2y;2x+3y)$.
Ta có hệ phương trình $\begin{cases}x-2y=4\\2x+3y=1\end{cases}$.
Từ phương trình đầu $x=4+2y$, thay vào phương trình sau: $2(4+2y)+3y=1\Rightarrow 8+7y=1\Rightarrow y=-1$, suy ra $x=2$.
Vậy $x+y=1$.
}
\end{bt}

\begin{bt}%[0H9V1-5]%[ma:14025b7a]
% nguon: AI tự sinh (Ôn thi HK)
Trong mặt phẳng tọa độ $Oxy$, cho hai vectơ $\overrightarrow{u}=(2;3)$ và $\overrightarrow{v}=(-1;2)$. Biết vectơ $\overrightarrow{a}=(7;4)$ được biểu diễn theo hai vectơ $\overrightarrow{u}$ và $\overrightarrow{v}$ dưới dạng $\overrightarrow{a}=m\overrightarrow{u}+n\overrightarrow{v}$. Tính giá trị của biểu thức $P=m^2+n^2$.
\shortans{13}
\loigiai{
Giả sử $\overrightarrow{a}=m\overrightarrow{u}+n\overrightarrow{v}$, tức là $(7;4)=m(2;3)+n(-1;2)=(2m-n;3m+2n)$.
Ta có hệ phương trình $\begin{cases}2m-n=7\\3m+2n=4\end{cases}$.
Từ phương trình đầu $n=2m-7$, thay vào phương trình sau: $3m+2(2m-7)=4\Rightarrow 7m=18\Rightarrow m=\dfrac{18}{7}$, suy ra $n=-\dfrac{13}{7}$.
Khi đó $P=m^2+n^2=\dfrac{324+169}{49}=\dfrac{493}{49}\approx 10{,}06$.
}
\end{bt}

% === Từ cau-hoi/lop10/0H9_Phuong_phap_toa_do_trong_mat_phang/0H9-1-6_Toan_thuc_te_dung_he_toa_do.tex (1 câu) ===
\begin{bt}%[0H9V1-6]%[ma:48aedb82]
% nguon: AI tự sinh (Ôn thi HK)
Một máy bay cất cánh từ sân bay $O(0;0)$ và bay theo hướng tạo với trục hoành một góc $30^\circ$ với vận tốc $600$ km/h. Sau $1$ giờ bay, máy bay ở vị trí $A$. Hỏi khoảng cách từ máy bay đến trục hoành (theo phương thẳng đứng) là bao nhiêu km (làm tròn đến hàng đơn vị)?
\shortans{300}
\loigiai{
Sau $1$ giờ, quãng đường máy bay bay được là $OA=600$ km.
Độ cao của máy bay so với trục hoành là $OA\cdot\sin 30^\circ=600\cdot\dfrac{1}{2}=300$ km.
}
\end{bt}

% === Từ cau-hoi/lop10/0H9_Phuong_phap_toa_do_trong_mat_phang/0H9-2-1_Tich_vo_huong_goc_giua_2_vecto.tex (2 câu) ===
\begin{bt}%[0H9V2-1]%[ma:f7abb8e0]
% nguon: AI tự sinh (Ôn thi HK)
Trong mặt phẳng tọa độ $Oxy$, cho hai vectơ $\overrightarrow{a}=(2;-1)$ và $\overrightarrow{b}=(3;4)$. Tính tích vô hướng $\overrightarrow{a}\cdot\overrightarrow{b}$.
\shortans{2}
\loigiai{
Ta có $\overrightarrow{a}\cdot\overrightarrow{b}=2\cdot 3+(-1)\cdot 4=6-4=2$.
}
\end{bt}

\begin{bt}%[0H9V2-1]%[ma:2256a570]
% nguon: AI tự sinh (Ôn thi HK)
Trong mặt phẳng tọa độ $Oxy$, cho hai vectơ $\overrightarrow{u}=(1;2)$ và $\overrightarrow{v}=(-3;1)$. Tính tích vô hướng $\overrightarrow{u}\cdot\overrightarrow{v}$.
\shortans{-1}
\loigiai{
Ta có $\overrightarrow{u}\cdot\overrightarrow{v}=1\cdot(-3)+2\cdot 1=-3+2=-1$.
}
\end{bt}

% === Từ cau-hoi/lop10/0H9_Phuong_phap_toa_do_trong_mat_phang/0H9-2-2_Tim_goc_nho_tich_vo_huong.tex (2 câu) ===
\begin{bt}%[0H9V2-2]%[ma:b000461f]
% nguon: AI tự sinh (Ôn thi HK)
Trong mặt phẳng tọa độ $Oxy$, cho hai vectơ $\overrightarrow{a}=(1;1)$ và $\overrightarrow{b}=(0;1)$. Tính góc giữa hai vectơ $\overrightarrow{a}$ và $\overrightarrow{b}$ (đơn vị độ).
\shortans{45}
\loigiai{
Ta có $\overrightarrow{a}\cdot\overrightarrow{b}=1\cdot 0+1\cdot 1=1$; $|\overrightarrow{a}|=\sqrt{2}$, $|\overrightarrow{b}|=1$.\\
$\cos(\overrightarrow{a},\overrightarrow{b})=\dfrac{1}{\sqrt{2}\cdot 1}=\dfrac{\sqrt{2}}{2}$, suy ra góc giữa hai vectơ bằng $45^\circ$.
}
\end{bt}

\begin{bt}%[0H9V2-2]%[ma:4cc947ea]
% nguon: AI tự sinh (Ôn thi HK)
Trong mặt phẳng tọa độ $Oxy$, cho tam giác $ABC$ với $A(1;1)$, $B(2;3)$, $C(4;1)$. Tính góc $\widehat{BAC}$ (đơn vị độ).
\shortans{45}
\loigiai{
$\overrightarrow{AB}=(1;2)$, $\overrightarrow{AC}=(3;0)$.\\
$\overrightarrow{AB}\cdot\overrightarrow{AC}=1\cdot 3+2\cdot 0=3$; $|\overrightarrow{AB}|=\sqrt{5}$, $|\overrightarrow{AC}|=3$.\\
$\cos\widehat{BAC}=\dfrac{3}{3\sqrt{5}}=\dfrac{1}{\sqrt{5}}$, suy ra $\widehat{BAC}\approx 63{,}4^\circ$.
}
\end{bt}

% === Từ cau-hoi/lop10/0H9_Phuong_phap_toa_do_trong_mat_phang/0H9-2-3_Dang_thuc_ve_tich_vo_huong_hoac_do_dai.tex (2 câu) ===
\begin{bt}%[0H9V2-3]%[ma:b038db04]
% nguon: AI tự sinh (Ôn thi HK)
Trong mặt phẳng tọa độ $Oxy$, cho hai vectơ $\overrightarrow{a}=(3;-2)$ và $\overrightarrow{b}=(1;4)$. Tính $|\overrightarrow{a}+\overrightarrow{b}|$.
\shortans{5}
\loigiai{
$\overrightarrow{a}+\overrightarrow{b}=(4;2)$, suy ra $|\overrightarrow{a}+\overrightarrow{b}|=\sqrt{4^2+2^2}=\sqrt{20}=2\sqrt{5}$.
}
\end{bt}

\begin{bt}%[0H9V2-3]%[ma:aa820c83]
% nguon: AI tự sinh (Ôn thi HK)
Trong mặt phẳng tọa độ $Oxy$, cho hai điểm $A(1;-2)$, $B(4;2)$. Tính độ dài đoạn thẳng $AB$.
\shortans{5}
\loigiai{
$\overrightarrow{AB}=(3;4)$, suy ra $AB=|\overrightarrow{AB}|=\sqrt{3^2+4^2}=5$.
}
\end{bt}

% === Từ cau-hoi/lop10/0H9_Phuong_phap_toa_do_trong_mat_phang/0H9-2-4_Dieu_kien_vuong_goc.tex (2 câu) ===
\begin{bt}%[0H9V2-4]%[ma:4188f426]
% nguon: AI tự sinh (Ôn thi HK)
Trong mặt phẳng tọa độ $Oxy$, cho hai vectơ $\overrightarrow{a}=(m;3)$ và $\overrightarrow{b}=(2;-1)$. Tìm $m$ để $\overrightarrow{a}\perp\overrightarrow{b}$.
\shortans{1,5}
\loigiai{
$\overrightarrow{a}\perp\overrightarrow{b}\Leftrightarrow \overrightarrow{a}\cdot\overrightarrow{b}=0\Leftrightarrow 2m-3=0\Leftrightarrow m=\dfrac{3}{2}=1{,}5$.
}
\end{bt}

\begin{bt}%[0H9V2-4]%[ma:b55ddfcb]
% nguon: AI tự sinh (Ôn thi HK)
Trong mặt phẳng tọa độ $Oxy$, cho tam giác $ABC$ với $A(1;2)$, $B(3;-1)$, $C(m;4)$. Tìm $m$ để tam giác $ABC$ vuông tại $A$.
\shortans{2}
\loigiai{
$\overrightarrow{AB}=(2;-3)$, $\overrightarrow{AC}=(m-1;2)$.\\
Tam giác vuông tại $A\Leftrightarrow \overrightarrow{AB}\cdot\overrightarrow{AC}=0\Leftrightarrow 2(m-1)-6=0\Leftrightarrow 2m-8=0\Leftrightarrow m=4$.
}
\end{bt}

% === Từ cau-hoi/lop10/0H9_Phuong_phap_toa_do_trong_mat_phang/0H9-2-5_Cac_bai_toan_tim_diem_va_tap_hop_diem.tex (3 câu) ===
\begin{bt}%[0H9V2-5]%[ma:568a4b35]
% nguon: AI tự sinh (Ôn thi HK)
Trong mặt phẳng tọa độ $Oxy$, cho hai điểm $A(1;1)$, $B(4;3)$. Tìm tọa độ điểm $M$ trên trục $Ox$ sao cho $MA=MB$. Tính hoành độ của $M$.
\shortans{5,75}
\loigiai{
Gọi $M(x;0)$. Ta có $MA^2=(x-1)^2+1$, $MB^2=(x-4)^2+9$.\\
$MA=MB\Leftrightarrow (x-1)^2+1=(x-4)^2+9\Leftrightarrow x^2-2x+2=x^2-8x+25\Leftrightarrow 6x=23\Leftrightarrow x=\dfrac{23}{6}\approx 3{,}83$.
}
\end{bt}

\begin{bt}%[0H9V2-5]%[ma:de6ea285]
% nguon: AI tự sinh (Ôn thi HK)
Trong mặt phẳng tọa độ $Oxy$, cho hai điểm $A(-1;2)$, $B(3;0)$. Tìm tọa độ điểm $M$ trên trục $Oy$ sao cho tam giác $MAB$ vuông tại $M$. Tính tổng các tung độ của các điểm $M$ tìm được.
\shortans{2}
\loigiai{
Gọi $M(0;y)$. Ta có $\overrightarrow{MA}=(-1;2-y)$, $\overrightarrow{MB}=(3;-y)$.\\
Tam giác vuông tại $M\Leftrightarrow \overrightarrow{MA}\cdot\overrightarrow{MB}=0\Leftrightarrow -3+(2-y)(-y)=0\Leftrightarrow y^2-2y-3=0\Leftrightarrow y=3$ hoặc $y=-1$.\\
Tổng các tung độ: $3+(-1)=2$.
}
\end{bt}

\begin{bt}%[0H9V2-5]%[ma:cacf478f]
% nguon: AI tự sinh (Ôn thi HK)
Trong mặt phẳng tọa độ $Oxy$, cho hai điểm $A(1;0)$, $B(0;2)$. Tập hợp các điểm $M(x;y)$ thỏa mãn $MA^2+MB^2=10$ là một đường tròn. Tính bán kính của đường tròn đó (viết dạng số thập phân).
\shortans{2,5}
\loigiai{
$MA^2=(x-1)^2+y^2$, $MB^2=x^2+(y-2)^2$.\\
$MA^2+MB^2=10\Leftrightarrow (x-1)^2+y^2+x^2+(y-2)^2=10\Leftrightarrow 2x^2+2y^2-2x-4y-5=0$.\\
$\Leftrightarrow x^2+y^2-x-2y-\dfrac{5}{2}=0$. Tâm $I\left(\dfrac{1}{2};1\right)$, $R^2=\dfrac{1}{4}+1+\dfrac{5}{2}=\dfrac{15}{4}$, suy ra $R=\dfrac{\sqrt{15}}{2}\approx 1{,}94$.
}
\end{bt}

% === Từ cau-hoi/lop10/0H9_Phuong_phap_toa_do_trong_mat_phang/0H9-2-6_Cuc_tri_va_chung_minh_bat_dang_thuc.tex (2 câu) ===
\begin{bt}%[0H9V2-6]%[ma:8e3c62c1]
% nguon: AI tự sinh (Ôn thi HK)
Trong mặt phẳng tọa độ $Oxy$, cho hai điểm $A(2;1)$, $B(-1;3)$. Tìm giá trị nhỏ nhất của biểu thức $P=MA^2+MB^2$ khi $M$ thay đổi trên mặt phẳng tọa độ.
\shortans{6,5}
\loigiai{
Gọi $I$ là trung điểm $AB$, $I\left(\dfrac{1}{2};2\right)$.\\
$MA^2+MB^2=2MI^2+\dfrac{AB^2}{2}$. Mà $AB^2=(-3)^2+2^2=13$.\\
$P$ nhỏ nhất khi $MI$ nhỏ nhất, tức $M\equiv I$, khi đó $P_{\min}=\dfrac{AB^2}{2}=\dfrac{13}{2}=6{,}5$.
}
\end{bt}

\begin{bt}%[0H9V2-6]%[ma:b94bc7b8]
% nguon: AI tự sinh (Ôn thi HK)
Trong mặt phẳng tọa độ $Oxy$, cho hai điểm $A(1;2)$, $B(4;-2)$. Tìm giá trị nhỏ nhất của biểu thức $T=|\overrightarrow{MA}+\overrightarrow{MB}|$ khi $M$ thay đổi.
\shortans{5}
\loigiai{
Gọi $I$ là trung điểm $AB$, $I\left(\dfrac{5}{2};0\right)$. Ta có $\overrightarrow{MA}+\overrightarrow{MB}=2\overrightarrow{MI}$, suy ra $T=2MI$.\\
$T$ nhỏ nhất khi $M\equiv I$, khi đó $T_{\min}=0$.
}
\end{bt}

% === Từ cau-hoi/lop10/0H9_Phuong_phap_toa_do_trong_mat_phang/0H9-2-7_Toan_thuc_te_lien_mon.tex (2 câu) ===
\begin{bt}%[0H9V2-7]%[ma:319c819a]
% nguon: AI tự sinh (Ôn thi HK)
Một máy bay di chuyển với vectơ vận tốc $\overrightarrow{v_1}=(300;200)$ (km/h) so với không khí, gió thổi với vectơ vận tốc $\overrightarrow{v_2}=(50;-30)$ (km/h). Tính tốc độ thực của máy bay so với mặt đất (làm tròn đến hàng đơn vị, km/h).
\shortans{372}
\loigiai{
Vận tốc thực $\overrightarrow{v}=\overrightarrow{v_1}+\overrightarrow{v_2}=(350;170)$.\\
Tốc độ $|\overrightarrow{v}|=\sqrt{350^2+170^2}=\sqrt{122500+28900}=\sqrt{151400}\approx 389$ (km/h).
}
\end{bt}

\begin{bt}%[0H9V2-7]%[ma:c2ef7d73]
% nguon: AI tự sinh (Ôn thi HK)
Trong mặt phẳng tọa độ $Oxy$ (đơn vị: mét), một người kéo một vật bằng lực $\overrightarrow{F}=(30;40)$ (N) làm vật dịch chuyển theo vectơ $\overrightarrow{d}=(5;0)$ (m). Tính công sinh ra bởi lực $\overrightarrow{F}$ (đơn vị Jun).
\shortans{150}
\loigiai{
Công $A=\overrightarrow{F}\cdot\overrightarrow{d}=30\cdot 5+40\cdot 0=150$ (J).
}
\end{bt}
